% Ad Familiares 6.20 --- headnote
% ad-familiares-06-20, end of July 45 BC, Astura

\textit{Cicero to C. Toranius, written from Astura at the
end of July (Quintilis) 45 BC --- Perseus dateline
\textit{Scr. Asturae ex. m. Quint. a. 709 (45)}. The
salutation \textit{CICERO TORANIO S.}\ identifies the
correspondent as Toranius, a Pompeian aedile of 64 (Cicero's
own consular year-of-canvass colleague) who, like the other
addressees of book 6, was waiting out the post-Pharsalus
settlement in exile. Astura is Cicero's seaside retreat south
of Antium, the small island-villa to which he had withdrawn
after the death of his daughter Tullia in February of this
year; nearly all of the \textit{De Finibus} and the
\textit{Tusculans} are being composed there. The letter is a
follow-up to one dispatched three days earlier through the
slaves of Cn.\ Plancius (\textit{Fam.}\ 6.21), and its
keynote is patience: stay where you are until Caesar's
return makes the next step clear.}

\textit{The letter divides into three movements --- a
practical case for not moving (\S 1), a brief summing-up
of the same point in a different register (\S 2), and a
philosophical close (\S 3) that turns the whole exchange
into a small \textit{consolatio}: ``whatever befalls us not
through our own fault we ought to bear bravely.'' ``Our
friend Cilo'' is a shared intermediary, otherwise
unidentified; the indirect manner in which Toranius's
situation is referred to (\textit{si recipiet ille se ad
tempus} --- ``if he comes back to himself in time'') is the
standard periphrasis of these months for Caesar's return
from Spain after Munda. The understated tetracolon
``\textit{te desiderant et diligunt et colunt}'' --- ``they
long for you, love you, cherish you'' --- closes the
letter on the affection of family, the one register the
political weather had not yet stripped from Cicero's
correspondents.}
